\begin{claim}{D.2.13}{Proved}{A finite obstruction to the canonical Hopf extension}
Put $B=K(\FF_2,2)$, $X=\Sigma B$, $Y=\Sigma^2B$,
$C_mX=\Omega^m\Sigma^mX$, and $S_m=\Sigma^2C_mX$.
Let $j:X\to QX$ be the stabilization unit and set
\[
 h_0=\eta_2\wedge\id_B:\Sigma^2X=\Sigma^3B\longrightarrow Y,
 \qquad \eta_2:S^3\longrightarrow S^2.
 \tag{D.2.13a}
\]
This map is stably essential, although its ordinary homology map is zero.
The suspension--loop counits give compatible extensions $h_m:S_m\to Y$
for $m=1,2$. The canonical $h_2$ extends through relative dimension
eleven in a relative CW approximation of $S_2\to\Sigma^2QX$,
but it does not extend across $S_2\to S_4$.

The coefficient calculation below would force the double adjunct of
any pointed extension $h:\Sigma^2QX\to Y$ of $h_0$ to fail to
preserve the Pontryagin square of the degree-five class $C$ specified
below. The additional multiplication-deviation argument of D.2.16
proves that no ordinary extension exists, already at $S_4$.
The general finite-detection argument remains useful: an extension
exists if and only if $h_0$ extends to every $S_m$, and a negative
answer is detected at a finite stage and on a finite relative subcomplex.
\end{claim}
\begin{proof}
Throughout the homology calculation coefficients are $\FF_2$;
write $\beta=\operatorname{Sq}_*^1$ and $T=\Omega^2Y$.
The argument uses Steenrod naturality for arbitrary maps. It uses
Pontryagin multiplicativity only for the specified double-loop map
$\Omega^2h_0$.

\medskip\noindent
\textbf{The Hopf attachment and its first extensions.}
The cofiber of $h_0$ is $\mathbb CP^2\wedge B$. Let $t$ be the
degree-two class of $\mathbb CP^2$, and let $u\in H^2B$ and
$v=\operatorname{Sq}^1u\in H^3B$. Serre's computation gives
$v\ne0$ and
\[
 \operatorname{Sq}^3(tu)=t^2v\ne0.
 \tag{D.2.13b}
\]
If the stable map were zero, its cofiber would split, compatibly with
the inclusion of $Y$, as $\Sigma^2\Sigma^\infty B\vee
\Sigma^4\Sigma^\infty B$. The unique degree-four class would then
have zero $\operatorname{Sq}^3$, since $\operatorname{Sq}^3u=0$.
This contradicts \textup{(D.2.13b)}. The homology map itself is zero
by the K\"unneth theorem applied to $\eta_2\wedge\id_B$.
The cohomology and operation conventions here are those of
\cite[\S1, nos.~2.1--2.4; \S2, nos.~8--9, pp.~199, 203--204]{I_Serre1953}.

For $0\le m\le2$ let
\[
 r_m=\Sigma^{2-m}\varepsilon^{(m)}:
 \Sigma^2\Omega^m\Sigma^mX\longrightarrow\Sigma^2X,
 \qquad h_m=h_0r_m.
 \tag{D.2.13c}
\]
The adjunction triangle identities show that these extend $h_0$;
naturality of the counit shows their compatibility under stabilization.
In particular the double adjunct of $h_2$ is
$g_2=\Omega^2h_0:C_2X\to T$.

\medskip\noindent
\textbf{The one prescribed label value.}
Use the following labels for classes of $\widetilde H_*B$, and for
their images in $H_*T$ under the unit:
\[
\begin{array}{c|rrrrrrr}
\text{label}&a&b&c&d&e&f&x\\
\text{degree}&2&3&4&5&5&6&6\\
\text{dual monomial}&u&v&u^2&w&uv&u^3&v^2
\end{array}
 \qquad w=\operatorname{Sq}^2v.
 \tag{D.2.13d}
\]
These are the complete positive-degree homology bases through degree
six. Serre's polynomial description and Cartan give
\[
 \beta b=a,\quad \beta c=\beta d=\beta e=\beta f=0,
 \quad\beta x=d+e,\quad \operatorname{Sq}_*^2c=a,
 \quad\overline\Delta c=a\otimes a.
 \tag{D.2.13e}
\]
For example $\operatorname{Sq}^1w=v^2=\operatorname{Sq}^1(uv)$,
which proves the formula for $\beta x$.

Write $\xi=\xi_1$ for the top operation on the homology of an
$E_2$-space, of degree $|\xi z|=2|z|+1$, and
$\{z,z'\}=\lambda_1(z,z')$ for its degree-one Browder bracket.
Let $g_0:X\to T$ be the double adjunct of $h_0$. Then
\[
 C=\sigma c\in H_5X,\qquad g_{0*}(C)=\xi a.
 \tag{D.2.13f}
\]
Here is a verification that does not identify reduced Hopf rotation
with an unreduced BV operation. The once-adjunct
$\varphi_0:\Sigma^2B\to\Omega\Sigma^2B$ has homology image
\[
 \varphi_{0*}(\sigma^2c)=(\sigma a)^2
       \in T(\widetilde H_*\Sigma B).
 \tag{D.2.13g}
\]
Indeed the quadratic part of a minimal Adams--Hilton attaching
differential is the reduced diagonal of the cone. In
$\mathbb CP^2\wedge B$, the product of the two old classes $tu$
is $t^2u^2$, giving precisely the displayed coefficient. The linear
part vanishes because $h_0$ is homology-zero, and a word of length
three has degree at least nine, whereas the input has degree six.
This is the relative suspension-cone form of the Adams--Hilton
calculation in \cite[\S3.1, p.~300; proof of Proposition~7.4,
p.~309]{FGHHalperinLemaire1995}. The tensor-algebra identification
is the James calculation \cite[\S3; Theorem~5.6 and Corollary~5.7]{I_James1955}.

Homology suspension $H_5T\to H_6\Omega Y$ sends $\xi a$ to
$(\sigma a)^2$. The free $E_2$ basis in degree five is
$d,e,ab,\xi a$. The first two have independent one-letter
suspensions, the last has the displayed two-letter suspension,
and $ab$ suspends to zero. Since $a,b$ are primitive,
$\overline\Delta(ab)=a\otimes b+b\otimes a\ne0$.
Thus the primitive kernel of homology suspension in this degree
is zero. The class $g_{0*}(C)$ is primitive because $X$ is a
suspension. This proves \textup{(D.2.13f)}.

\medskip\noindent
\textbf{Two coefficient identities in the target.}
The free $E_2$ homology theorem describes $H_*T$ as the polynomial
algebra on a basis of the free restricted, degree-one Lie algebra
on the labels of $\widetilde H_*B$; use
\cite[Chapter III, Theorems~3.1--3.2 and Lemma~3.8]{I_CohenLadaMay1976}.
In particular pure restricted powers of a label are independent
of the Lie monomials containing a bracket. The exact top-operation
relations, including their Browder terms, are
\[
 \beta\xi z=(|z|\bmod2)z^2+\{\beta z,z\},\qquad
 \operatorname{Sq}_*^2\xi z=\xi(\beta z)+\{z,\operatorname{Sq}_*^2z\}.
 \tag{D.2.13h}
\]
The bracket has the Steenrod Cartan formulas and $\{z,z\}=0$.
These are \cite[Chapter III, Theorems~1.2--1.4,
pp.~216--219]{I_CohenLadaMay1976}; their homology-suspension formula
also supplies the operation used above.

Let $\lambda:H_{10}T\to\FF_2$ select the coefficient of
$(\xi a)^2$, and let $\mu:H_{11}T\to\FF_2$ select the coefficient
of $\xi(\xi a)$. We claim
\[
 \lambda\beta=\mu,\qquad
 \mu\operatorname{Sq}_*^2=0:H_{13}T\longrightarrow\FF_2.
 \tag{D.2.13i}
\]
All basis possibilities can be checked in these degrees as follows.
The degree-five Lie generators are $d,e,\xi a$; hence the pure
top generators in degree eleven are $\xi d,\xi e,\xi(\xi a)$.
Their Bocksteins are respectively $d^2,e^2,(\xi a)^2$ by
\textup{(D.2.13e),(D.2.13h)}. Bocksteins of bracket expressions
retain a bracket factor after the Cartan formula is expanded, so
cannot contribute the specified pure monomial. A decomposable
class could contribute only through a degree-six Bockstein
preimage of $\xi a$. There is none: the complete degree-six
basis is
\[
 f, x, ac, b^2, a^3, \{a,b\},
 \tag{D.2.13j}
\]
whose Bocksteins are respectively $0,d+e,0,0,0,0$.
This proves the first identity.

For the second, products have decomposable Steenrod images,
labels remain labels, and brackets remain bracket expressions.
A pure top generator of degree thirteen has the form $\xi z$
with $z$ a degree-six Lie generator, hence $z=f,x,$ or $\{a,b\}$.
Its square-operation image modulo the bracket ideal is
$\xi(\beta z)$. The displayed degree-six list shows that
$\beta z$ has no $\xi a$ component; it is a combination of
$d,e$. Thus $\xi(\beta z)$ has no $\xi(\xi a)$ component.
This proves the second identity, and in particular
$\lambda\beta\operatorname{Sq}_*^2=0$ on $H_{13}T$.

\medskip\noindent
\textbf{The obstruction already at the fourth stage.}
Use $C$ also for its image in $H_5C_mX$, and put
$A_0=\sigma a\in H_3C_mX$. They satisfy
$\beta C=\beta A_0=0$ and $\operatorname{Sq}_*^2C=A_0$.
At the fourth stage $\xi_3C=Q^8C$ is the \emph{top} operation.
Its exact Nishida formula is
\[
 \operatorname{Sq}_*^2\xi_3C
       =Q^6C+\lambda_3(C,A_0),\qquad
 \beta Q^6C=C^2.
 \tag{D.2.13k}
\]
The omitted possible ordinary term is $Q^7\beta C=0$.
The coefficient of $Q^6C$ is odd. These are the top and non-top
formulas of \cite[Chapter III, Theorems~1.1--1.3,
especially pp.~217--218]{I_CohenLadaMay1976}.
The Bockstein of the Browder correction is zero because both
its inputs are Bockstein-zero. Consequently
\[
 \beta\operatorname{Sq}_*^2\xi_3C=C^2
       \quad\hbox{in }H_{10}C_4X.
 \tag{D.2.13l}
\]
This step retains the top-operation correction; the uncorrected
Nishida formula alone would not justify a fourth-stage assertion.

Suppose $h_2$ extended to $S_4$. Taking double adjuncts would give
an ordinary pointed map $G:C_4X\to T$ restricting to
$g_2=\Omega^2h_0$. This specified restriction is multiplicative,
so \textup{(D.2.13f)} implies
$\lambda G_*(C^2)=\lambda(\xi a)^2=1$.
But Steenrod naturality and \textup{(D.2.13i),(D.2.13l)} give
\[
 \lambda G_*(C^2)
   =\lambda\beta\operatorname{Sq}_*^2G_*(\xi_3C)=0,
 \tag{D.2.13m}
\]
a contradiction. No multiplicativity was assumed for $G$.
The relation detecting this obstruction lies in homology degree
thirteen before the two suspensions.

In the full infinite-loop source the ordinary Nishida identities
give $\operatorname{Sq}_*^2Q^8C=Q^6C$ and $\beta Q^6C=C^2$.
Thus every ordinary extension $g:QX\to T$ would satisfy
\[
 \lambda g_*(C^2)=0,
 \qquad \lambda\bigl(g_*(C)^2\bigr)=1.
 \tag{D.2.13n}
\]
In particular it cannot be an H-map. For an arbitrary map the
first value is not constrained to equal the second; this coefficient
calculation alone therefore does not decide the full extension
question. D.2.16 uses an actual multiplication-deviation map and
additional cohomology identities to exclude every ordinary extension.

\medskip\noindent
\textbf{The positive range through relative dimension eleven.}
Freudenthal makes $B\to\Omega\Sigma B$ three-connected and
$\Omega\Sigma B\to\Omega^2\Sigma^2B$ four-connected.
Their composite is three-connected. Since $\pi_3B=0$, it follows
that $\pi_5Y=\pi_3T=0$; see
\cite[Corollary~4.24]{FGHHatcher2002} for the connectivity form.
Choose a bottom cell $i:S^3\to X$ and set $X'=X/S^3$.
The quotient is three-connected, since $i$ surjects onto
$H_3X=\ZZ/2$ and the quotient is simply connected.
The map $h_0\Sigma^2i$ is null, so the cofiber property gives
$\bar h:\Sigma^2X'\to Y$ with $\bar h\Sigma^2q\simeq h_0$,
where $q:X\to X'$.

The stabilization
$t':\Sigma^2\Omega^2\Sigma^2X'\to\Sigma^2QX'$
is eleven-connected: $\Sigma^2X'$ is five-connected,
Freudenthal gives eleven-connectivity before looping, and
looping twice and suspending twice restore that same bound.
The homotopy pullback
\[
 P=(\Sigma^2QX)\times^h_{\Sigma^2QX'}
                 (\Sigma^2\Omega^2\Sigma^2X')
 \tag{D.2.13o}
\]
therefore maps eleven-connectedly to $\Sigma^2QX$.
Naturality supplies a lift of $S_2$ to $P$. The other projection,
followed by the counit and $\bar h$, gives $P\to Y$ whose
restriction is $h_2$, by counit naturality and
$\bar h\Sigma^2q\simeq h_0$.
Replace $P\to\Sigma^2QX$ by a fibration. Its fiber is
ten-connected, so the specified lift extends over every relative
cell of dimension at most eleven. Composition with $P\to Y$
gives the asserted positive range. This argument uses an actual
pullback; it does not suspend a fiber sequence as a fiber sequence.

\medskip\noindent
\textbf{Finite detection for arbitrary choices of extensions.}
Serre's class theorem gives finite $2$-primary homology for $B$
in every positive degree and finite $2$-primary homotopy groups
for the simply connected space $Y$. The exact interfaces are
\cite[Chapter I, \S6.4; Chapter II, Proposition~8; Chapter III,
Theorem~1 and Corollaries~1--2, pp.~265, 271--274]{I_SerreClasses1953}.
All the sources above admit countable CW models. For a finite
CW complex $K$, Postnikov obstruction theory makes $[K,Y]_*$
finite and makes $\pi_1\Map_*(K,Y)$ finite at every component.
For an exhaustion $K=\bigcup K_s$ by finite subcomplexes,
the mapping-tower argument consequently gives
\[
 [K,Y]_*\cong\lim_s[K_s,Y]_*.
 \tag{D.2.13p}
\]
The obstruction to injectivity is a $\lim^1$ of finite groups;
its tower is Mittag--Leffler, so it vanishes. Surjectivity follows
by realizing compatible maps and homotopies on the telescope.
The right side gives $[K,Y]_*$ a compact Hausdorff profinite
topology. Restriction along any map of countable CW sources is
continuous, since its restriction to a finite subcomplex factors
through a finite subcomplex after cellular approximation.

Let $E_m\subseteq[S_m,Y]_*$ consist of the classes restricting
to $[h_0]$ on $S_0=\Sigma^2X$. These are closed compact sets, with
continuous restriction maps $E_{m+1}\to E_m$. If every $E_m$
is nonempty, the compatibility conditions in $\prod_mE_m$ have
the finite intersection property: any finite list is satisfied
by choosing a class at its largest index and restricting it.
Hence $\lim_mE_m$ is nonempty, without a surjectivity assumption
on the bonding maps. The mapping-telescope surjection realizes
such a compatible family on
$\operatorname{hocolim}_mS_m\simeq\Sigma^2QX$.
Conversely an extension clearly restricts to all finite stages.
At a fixed stage, apply the same finite Postnikov and compactness
argument to an exhaustion of a mapping-cylinder CW pair for
$S_0\to S_m$. Failure to extend is detected on a finite subpair.
This equivalence concerns \emph{all} choices in $E_m$. The
obstructed canonical choice above alone does not prove $E_4$ empty;
that stronger conclusion is proved separately in D.2.16.
\end{proof}
